\chapter{Anatomical structure, initialisation, and learned correspondence}
\label{ch:structure}

Chapter~\ref{ch:pipeline} described the pipeline as four sequential stages.
This chapter reports the technical work behind the first three, organised
around a single question asked with escalating power: \emph{where can the
fitter be given more of the information the surface objective lacks?} First
as fixed anatomical priors; then as a better starting point; then inside the
optimiser itself; then as a descriptor learned from data; and finally as a
counterfactual test of how much each source of information is actually
worth. Every result below was tested against a controlled benchmark before
being adopted or rejected; several are negative results, reported because a
negative result under a controlled comparison rules out a hypothesis as
usefully as a positive one confirms it.

\subsection*{Anatomical priors}
\section{Chain-aware anatomical assignment}
\label{sec:chain-aware}

Several downstream analyses in this report (the per-part correspondence
audit, the multi-level evaluation suite in \S\ref{sec:eval-suite})
depend on a prior step: deciding which anatomical structure a given surface
point belongs to. A naïve nearest-segment rule, which assigns each point to
whichever skeleton bone is geometrically closest, fails systematically on an
articulated body plan: a point on a middle leg can sit closer in raw
distance to the wrong leg, or to the body, than to its own leg, because the
rule treats every bone independently rather than as part of a connected
kinematic chain. A \textbf{chain-aware} rule, which resolves ambiguous
points by their position along the connected chain rather than by raw
distance, was tested against it on the same labelled template and adopted:
macro-averaged F1 improves from $0.656$ to $0.831$, concentrated exactly
where the naïve rule was weakest (middle-leg recall $0.349\to0.768$,
body/core recall $0.366\to0.663$). The lesson generalises beyond this one
rule: an articulated body plan is not a set of independent parts, and
geometric rules that treat it as one will fail systematically at the joints
between parts.

\begin{algbox}{2: Chain-aware anatomical assignment}
\begin{enumerate}[nosep,leftmargin=1.4em]
\item For each surface point $p$, compute distance to every bone segment.
\item Group segments into kinematic chains (each leg, antenna, mandible, body axis).
\item Assign $p$ to the chain whose segments are collectively closest, using
the chain's connected extent rather than the single nearest segment.
\item Within that chain, assign to the nearest segment.
\end{enumerate}
\end{algbox}

\section{The Shape Diameter Function as an anatomical prior}

A separate line of investigation asked whether a local-thickness descriptor could
support anatomical assignment more robustly than surface-based rules. The
\term{shape diameter function}~\cite{sdf} assigns every surface point a
value related to the local thickness of the solid behind it: the value $\mathrm{SDF}(x)$ is obtained by casting rays into the solid
behind a surface point and taking a robust average of chord lengths, so thin
structures (legs, antennae) take small values, the bulky body takes large
ones, independently of pose; and, unlike the model's own vertex
correspondence, it is computable directly on a raw, unregistered scan. It
proved reliable for one purpose and unreliable for another. It reproduces
well across real scans (correlation $r\approx0.77$--$0.94$, $87$--$95\%$
agreement, under $1.3$\,s per specimen) and separates body from appendage on
the template with high accuracy (AUC $\approx0.973$, balanced accuracy
$\approx95.1\%$). But thresholding the field to separate \emph{individual}
appendage instances (one leg from its neighbour, left mandible from right)
fails on real scans: only $2$ of $12$ specimens had all four
expert-identified landmarks fall into distinct connected components, with
the two mandibles merged into one component in $9$ of $12$. The descriptor
is retained as a reliable, cheap semantic signal for \emph{which class of
structure} a region belongs to; it is not a mechanism for separating
individual instances of a repeated structure, because connected components cannot separate two
structures that are physically fused in the scan (laterality is not encoded in local
thickness). This is the instance-level correspondence problem this report's other sections address by different
means.

\section{Model configuration: joint limits and the shape prior}
\label{sec:model-configuration}

The production template used throughout this report is a joint-limited
SMIL model: $10{,}235$ vertices, $55$ joints, a $25$-dimensional shape space,
with an accompanying joint regressor, kinematic tree, skinning weights,
bilateral symmetry mapping, and, as the specific addition evaluated here,
authored per-joint rotation limits. Consuming these limits as a hard prior
during optimisation is the rotation-limit intervention already reported as
adopted in \S\ref{sec:eval-suite} and analysed mechanistically in \F{F1};
this section reports two related configuration choices tested alongside it.

\textbf{Joint limits and bilateral symmetry improve plausibility, not
correspondence.} Constraining joint rotation to anatomically authored ranges
measurably improves physical plausibility (fewer anatomically impossible
poses), and a companion bilateral (midline) symmetry constraint measurably
improves left/right consistency. Neither, on its own, resolves the
underlying correspondence problem characterised in \F{F1}, \F{F3}, and
\F{F6}: a pose can be entirely plausible, and entirely symmetric, while
still assigning the wrong vertex to the wrong anatomical location. This is
the same distinction as \F{F1} applied to a different pair of constraints:
plausible articulation is not the same property as correct articulation.

\textbf{A shape prior was added, then removed.} A Gaussian penalty on the
shape parameters $\bm\beta$ was tested as a stabiliser, on the standard
assumption that penalising extreme shape values improves numerical
robustness. It was subsequently removed after being found to suppress
genuine, biologically meaningful shape variation, the same variation
Chapter~\ref{ch:validation} depends on to distinguish specimens from one
another at all. This is reported as a deliberate reversal rather than an
oversight: a regularisation term can improve an optimisation diagnostic
while actively damaging the property the model exists to preserve, and for
a morphology-measurement pipeline specifically, that trade-off is not worth
making.


\subsection*{A better starting point}
\section{A learned pose initialiser, and where it breaks}
\label{sec:pointnet-init}

\S\ref{sec:init-stage} introduced a \term{PointNet}-style~\cite{pointnet}
initialiser: a point-cloud encoder that reads the raw target points,
extracts a permutation-invariant feature per point, pools them into one
global descriptor, and regresses an initial pose and shape directly from
it, replacing a fixed canonical starting pose with a learned, input-aware
one. The motivation follows from \F{F4}: for the distal chain, initialisation determines
which target surface region becomes each vertex's first attractor, so a
better starting point should, in principle, reduce how often the optimiser
locks onto the wrong one.

\begin{figure}[!htb]
\centering
\begin{tikzpicture}[every node/.style={font=\scriptsize},
 box/.style={draw,rounded corners=2pt,minimum height=8mm,align=center,fill=white},
 arr/.style={-{Stealth[length=1.6mm]}}]
\node[box] (a) {point cloud\\$N\times3$};
\node[box,right=6mm of a] (b) {shared MLP\\per point};
\node[box,right=6mm of b] (c) {symmetric\\max-pool};
\node[box,right=6mm of c] (d) {global\\descriptor};
\node[box,right=6mm of d] (e) {regression head\\$\hat{\bm\theta}$};
\node[box,right=6mm of e] (f) {SMILify\\optimiser};
\foreach \x/\y in {a/b,b/c,c/d,d/e,e/f}{\draw[arr] (\x)--(\y);}
\end{tikzpicture}
\caption{Learned initialiser. Max-pooling makes the descriptor invariant to
point ordering, which a raw scan does not have.}
\label{fig:pointnet}
\end{figure}


\begin{investigation}{I11}{Does the learned initialiser improve registration,
and does a synthetic improvement transfer to real specimens?}
\invQuestion{Under controlled, exact synthetic ground truth, does the
learned initialiser outperform the existing default; and does that
advantage hold on real scans?}
\invMethod{Synthetic benchmark: learned vs.\ default initialisation, exact
ground truth known. Real-data benchmark: 50 real specimens, learned vs.\
zero (default) initialisation at three levels of synthetic pose
perturbation used during training ($\approx\!23^\circ$, then retrained at
$\approx\!13.8^\circ$ and $\approx\!8.85^\circ$).}
\invResult{On synthetic data the learned initialiser wins clearly: leg
accuracy improves from $0.857$ to $0.901$ and Chamfer distance from
$0.000264$ to $0.000226$. On the 50-specimen real benchmark at the
$23^\circ$ training regime, this reverses: learned initialisation
($0.000697$ mean Chamfer) is \emph{worse} than the zero-initialised default
($0.000637$), winning only 20 of 50 specimens. Retraining at smaller
perturbation magnitudes moves the comparison toward parity rather than
robust superiority.}
\invVerdict{REJECT}{The observed failure is consistent with a
synthetic-to-real distribution mismatch: performance did not transfer from
the tested synthetic training regime to real scans, and the
perturbation-magnitude sweep supports that interpretation. The experiments
do not establish whether the architecture itself would generalise under a
more representative training distribution, only that it did not under the
one tested. Not deployed in production as evaluated. The general lesson is
retained as a standing caution: a better initialiser can, in principle,
improve optimisation without addressing correspondence at all, and a
synthetic evaluation cannot certify real-world transfer by itself.}
\end{investigation}

A controlled follow-up sweep, holding total pose-error magnitude fixed and
varying only where along the kinematic chain it was placed, is what produced
the proximal/distal result reported as \F{F4} in \S\ref{sec:articulation}:
the same error placed proximally degrades leg accuracy far more (to $0.768$)
than the identical error placed distally ($0.939$), because a proximal
mistake propagates through every joint downstream of it.


\subsection*{A better-behaved optimiser}
\section{Sampling starvation and robust optimisation}
\label{sec:sampling}

\textbf{Sampling starvation.} The surface term is a Monte Carlo estimate,
$\mathcal{L}_{\mathrm{ch}}\approx\frac1N\sum_i d(p_i,T)^2$ with
$p_i\sim p(x)$. If $p$ is proportional to surface area, the expected loss
$\mathbb{E}_{x\sim p}[\mathrm{err}(x)]$ weights the body and head far more
than distal legs or antennae, so a tiny structure can be almost invisible to
the optimiser: a second route to the under-representation of \F{F5}, this
time in the sampling scheme rather than the mesh. Reserving a fixed quota of
the global budget for distal points improved distal legs but starved
antennae and head, which shared that budget; restricting the quota to within
the leg budget avoided the trade-off, yet did not clearly improve on simply
splitting distal sampling. The conclusion is narrow: starvation is a real
confound, but controlling it did not create anatomical correspondence.

\textbf{Robust optimisation.} A related question is how the objective
should respond to points that are simply wrong (occluded regions,
scanning artefacts, disconnected debris) rather than merely
under-sampled. \term{Graduated non-convexity} (GNC)~\cite{gnc} is a robust-optimisation
strategy that treats every point correspondence as suspect initially and
progressively hardens its confidence in inliers as optimisation proceeds,
rather than trusting every nearest-point match from the first iteration.
Applied globally across the whole body, it produced undesirable behaviour
specifically around the antennae; restricted to the leg point set alone
(\emph{leg-only, topology-free}), it improved robustness where applied,
under synthetic corruption benchmarks with 30\% and 60\% of points dropped:
accuracy rose from $0.602$ to $0.797$ (Drop30) and from $0.411$ to $0.607$
(Drop60). Localising the robust treatment to the anatomical subsystem where
corrupted points actually occur outperformed applying it uniformly; the
result is retained as a promising candidate for leg registration
specifically, not adopted as a general replacement for point-matching across
the whole body.


\subsection*{Correspondence learned from data}
\section{From geometric to learned correspondence}
\label{sec:learned-correspondence}

Every correspondence mechanism considered so far in this report is
geometric: a vertex is matched to whichever target point is nearest to it,
possibly after the chain-aware and shape-diameter refinements of the previous
sections. A separate line of work asked whether a \emph{learned}
descriptor, a small network trained to produce a feature vector per
point such that corresponding anatomical locations produce similar vectors
could do better, and pursued that question through four increasingly
demanding tests.

The first attempt, a handcrafted baseline rather than a learned one,
normalised each point's coordinates within its own anatomical part (by that
part's local length and radius) before matching, on the hypothesis that
canonicalising away each part's specific proportions would make
correspondence easier to learn. It performed \emph{worse} than raw
coordinates: normalising away a part's specific geometry destroys
information that, it turns out, itself carries correspondence signal. The
first genuinely learned descriptor, trained with an off-the-shelf objective,
then failed when deployed, for a mundane but instructive reason: a mismatch
between the density of points it was trained on and the density it was
evaluated and deployed on, not a failure of the learning approach itself.
This motivated a redesigned, within-part training objective, tested next.

Two tasks must be kept apart. \emph{Pairwise retrieval} matches each query
point to its nearest neighbour in descriptor space,
$q_i\mapsto\arg\min_j\|f(q_i)-f(t_j)\|$. \emph{Dense correspondence} needs a
whole assignment $\{q_i\}_{i=1}^N\mapsto\{t_{\pi(i)}\}_{i=1}^N$ that is
globally consistent and covers the surface. Good retrieval does not
guarantee a consistent assignment.

\begin{algbox}{3: Learned-correspondence evaluation}
\begin{enumerate}[nosep,leftmargin=1.4em]
\item Compute descriptors $f$ for template and target points.
\item Retrieve candidates by nearest neighbour in descriptor space.
\item Score pairwise accuracy at a joint (retrieval).
\item Deploy the votes as a dense field over the whole mesh; measure coverage
and error against the geometric baseline across all specimens.
\end{enumerate}
\end{algbox}

\begin{investigation}{I14}{Can a descriptor trained specifically within one
anatomical part outperform raw local geometry at correspondence within that
part?}
\invQuestion{Trained and evaluated with an identical backbone and warm-start
procedure to the earlier attempt, does a within-part learned descriptor
improve correspondence accuracy at a specific, hard joint?}
\invMethod{Held-out coxa correspondence error, within-part learned
descriptor vs.\ the earlier learned baseline, same backbone, bootstrap
resampling for a confidence interval.}
\invResult{Coxa correspondence error falls from $0.2689$ to $0.2247$, a
$16.4\%$ improvement (bootstrap 95\% CI $[12.6\%, 19.8\%]$,
$p=7.8\times10^{-26}$), and the same descriptor outperforms the earlier
baseline across other segments too.}
\invVerdict{HOLD}{A learned local descriptor extracts correspondence
information that raw local geometry alone does not provide. This is
retained as evidence that learning has something to contribute, not as a
deployed mechanism: \S\ref{sec:learned-correspondence} next tests exactly
this descriptor as a dense correspondence field and rejects it at that
level. This is the strongest positive evidence in the investigation that
learning captures signal here, and, as the next result shows, it is also
where the harder problem begins.}
\end{investigation}

\begin{figure}[!htb]
\centering
\begin{tikzpicture}[every node/.style={font=\scriptsize},
 n/.style={draw,rounded corners=2pt,minimum height=7mm,align=center,fill=white,inner sep=3pt}]
\node[n,fill=green!12] (a) {Local test\\checkmark};
\node[n,below=3mm of a] (a2) {within-part descriptor};
\node[n,below=3mm of a2] (a3) {better pairwise retrieval};
\node[n,below=3mm of a3,fill=green!12] (a4) {$+16.4\%$ coxa};
\node[n,right=22mm of a,fill=red!10] (b) {Deployment\\$\times$};
\node[n,below=3mm of b] (b2) {dense correspondence};
\node[n,below=3mm of b2] (b3) {votes collapse};
\node[n,below=3mm of b3,fill=red!10] (b4) {$28.1\%\to1.8\%$ coverage};
\node[below=4mm of $(a4)!0.5!(b4)$,font=\small\bfseries] {local success $\neq$ global consistency};
\end{tikzpicture}
\caption{The same descriptor, evaluated at two levels of abstraction. Strong
local (pairwise) retrieval does not imply a usable dense correspondence
field.}
\label{fig:local-vs-global}
\end{figure}

\begin{investigation}{I15}{Does strong pairwise retrieval accuracy translate
into a usable, dense correspondence field at deployment?}
\invQuestion{Deployed as a dense, whole-mesh correspondence mechanism rather
than evaluated pairwise at one joint, does the within-part descriptor still
outperform the existing geometric baseline?}
\invMethod{Full-mesh deployment benchmark across 48 specimens, learned
descriptor vs.\ geometric nearest-point baseline; a follow-up analysis
isolating how much of any gap is attributable to differences in the
candidate point set each method was allowed to match against, versus the
descriptor's votes themselves.}
\invResult{The learned descriptor performs \emph{worse} at deployment scale.
Candidate-set differences account for only about $2.9\%$ of the observed
gap, against an overall gap of roughly $145\%$; at the coxa specifically,
surface coverage collapses from $28.1\%$ under the geometric baseline to
$1.8\%$ under the learned descriptor.}
\invVerdict{REJECT}{The learned votes collapse onto a small subset of
points instead of producing a dense, globally consistent field across the
whole mesh. Pairwise retrieval accuracy at a single joint (Investigation~I14) is not
the same property as dense, globally consistent correspondence across an
entire specimen, and the two must be evaluated separately rather than
assumed to track each other.}
\end{investigation}

Read together, these four steps identify a specific, narrow gap rather than
a blanket failure of learning: local, within-part descriptors carry real
information (Investigation~I14) that geometric nearest-matching does not,
but turning that local signal into a dense, mesh-wide correspondence
field (global consistency, not just local accuracy) is an unsolved
engineering problem in its own right (Investigation~I15), and is the
concrete direction this report's outlook (Chapter~\ref{ch:conclusion})
recommends pursuing further.


\subsection*{Counterfactual diagnosis}
\section{Counterfactual bottleneck analysis}
\label{sec:oracle}

Every experiment so far asks \emph{can an intervention make the system
better?} An \term{oracle intervention} asks a different, more powerful
question: \emph{if this one piece of information were solved exactly, how
much of the failure would remain?} The fitter is supplied, at inference,
with the correct answer to exactly one sub-question (information it could
not obtain on real data), and the residual error localises the bottleneck
causally, without proposing a deployable fix. Two oracles were run on the
synthetic corpus, where ground truth is exact, escalating in how much
information is supplied.

\begin{center}
\small
\begin{tabular}{@{}p{0.17\linewidth}p{0.20\linewidth}p{0.30\linewidth}p{0.20\linewidth}@{}}
\toprule
\textbf{Oracle} & \textbf{Information supplied} & \textbf{What remains} & \textbf{Evidence} \\
\midrule
Regional identity & ``which structure is this?'' & Localisation error inside
  the correct structure & Leg accuracy $0.867\to0.937$ (\F{F3}) \\
Dense correspondence & ``where exactly, for every vertex?'' & Objective,
  deformation and trajectory effects, even with correspondence solved &
  Residual error persists \\
Expert joints (\S\ref{sec:skeleton-underdetermined}) & ``where is the true
  skeleton?'' & Incomplete propagation to unsupervised regions &
  $25.1\%\to17.4\%$ WL, leave-one-region-out null (\F{F6}) \\
\bottomrule
\end{tabular}
\end{center}

\textbf{Regional identity.} Every target point is labelled with the region it
truly belongs to (which leg, which body segment), leaving the optimiser to
decide only where \emph{within} that region the point sits. Leg accuracy
improves from $0.867$ to $0.937$: a real but partial recovery. A large share
of correspondence error survives even when ``which structure is this?'' is
answered perfectly, so the harder remaining question is \emph{where within
the structure} a point belongs: regional identity and within-region
localisation are distinct problems (\F{F3}).

\textbf{Dense correspondence.} A second oracle supplies near-exact per-vertex
correspondence, including at very high weight against the rest of the
objective. Residual error remains here too: the objective's other terms, the
deformation freedom of \S\ref{sec:deformation-freedom}, and the optimisation
trajectory continue to limit the fit even when correspondence is no longer
the constraint. This rules out the simplest reading of this report's central
theme, that correspondence is the \emph{only} missing ingredient, and
locates the problem in an interaction between correspondence, objective,
trajectory and deformation. A third, real-data oracle along the same logic
(supplying expert joint locations) is reported in
\S\ref{sec:skeleton-underdetermined}, where the same escalating pattern holds:
more information supplied recovers more accuracy, but does not propagate
beyond where it is supplied.

